% ============================================================
%  参考文献
% ============================================================
\begin{thebibliography}{99}

\bibitem{ref1}
作者一, 作者二.
\emph{论文标题} [J].
期刊名, 年份, 卷(期): 起止页码.

\bibitem{ref2}
作者.
\emph{书名} [M].
出版地: 出版社, 年份.

\bibitem{ref3}
作者.
\emph{论文标题} [C].
会议名称, 会议地点, 年份.

% 若使用 AI 工具，请在正文相应位置标注，并按竞赛要求在参考文献中列出：
% [编号] 工具名称，版本/型号，开发机构/公司，使用日期。
%
% 示例：
% \bibitem{ai-tool}
% DeepSeek, DeepSeek-R1-0528, 深度求索（DeepSeek）, 2025-09-05.
%
% 更多参考文献按需添加
% 使用 \cite{ref1} 在正文中引用

\end{thebibliography}

% 若未使用 AI 工具，保留下方声明；若已使用 AI 工具，请删除本声明，
% 并在支撑材料中另附 PDF 文件“AI 工具使用详情”。
\paragraph*{人工智能工具使用声明}
本参赛队未使用任何 AI 工具。
